\documentclass[11pt, a4paper]{report}

\usepackage[utf8]{inputenc}
\usepackage[T1]{fontenc}
\usepackage[french]{babel}

% Appel des packages dans le dossier preambule
\def\packagepath{../../../../../preambule}   % path du package princial

\usepackage{\packagepath/page_layout} % <--- Nouveau
\usepackage{\packagepath/config_info}
\usepackage{\packagepath/cadres_cours}
\usepackage{\packagepath/zones_reponse}
\usepackage{\packagepath/config_ds}
\usepackage{\packagepath/config_maths}

% --- PILOTAGE ---
%\def\visibleornot{visible}    % visible or invisible
\ifdefined\visibleornot\else\def\visibleornot{visible}\fi


\def\documentpath{./Sources_Latex}
\graphicspath{{./Images/}}


\def\ptsexoA{0}



\begin{document}

\def\level{Premiere}  % Classe à droite
\def\course{NSI}
\def\chaptername{Representation des entiers relatifs} % Titre au centre
\def\docname{RD03~--~TD 1}        % Identifiant à gauche

\pagestyle{DS_LP}

\NomPrenom{}

\bigskip

%%=============================================================
\emph{Sauf mention contraire, les mots binaires seront toujours des mots de 8 bits et représenteront des entiers relatifs codés en complément à 2.}

\medskip

\begin{minipage}{0.48\linewidth}

        \begin{exercice_cours}[Conversion Décimal -> Binaire]

        \begin{tabular}{p{3.6cm}c}
    \hline
    Nombre entier & Nombre binaire \\
    \hline
    \verb|-37| & \reponseLigne[2.5cm]{\texttt{11011101}} \\
    \verb|-132| & \reponseLigne[2.5cm]{\texttt{10110100}} \\
    \verb|-27| & \reponseLigne[2.5cm]{\texttt{11100101}} \\
     \hline
        \end{tabular}
        \end{exercice_cours}
    \end{minipage}\hfill
    \begin{minipage}{0.48\linewidth}
                \begin{exercice_cours}[Conversion Binaire -> Décimal]
        \begin{tabular}{p{3.6cm}c}
    \hline
    Nombre binaire & Nombre entier \\
    \hline
    \verb|10011101| & \reponseLigne[2.5cm]{\texttt{-37}} \\
    \verb|10110110| & \reponseLigne[2.5cm]{\texttt{-132}} \\
    \verb|11101001| & \reponseLigne[2.5cm]{\texttt{-23}} \\
    \hline
        \end{tabular}
                \end{exercice_cours}
    \end{minipage}
%%=============================================================

\bigskip



\begin{exercice_cours}[Répondre aux questions suivantes:]

    \begin{itemize}
        \item combien de valeurs différentes peut-on représenter avec 2 bits ? Avec 6 bits ? Avec 8 bits ? Avec 32 bits ? 
        \reponseLigne[0.8\linewidth]{\texttt{ respectivement 4, 64, 256 et 4294967296}}
        \item quel est le plus petit entier relatif qu'on peut représenter avec 8 bits ? \hfill\reponseLigne[2cm]{\texttt{-128}}
        \item quel est le plus grand entier relatif qu'on peut représenter avec 8 bits ? \hfill\reponseLigne[2cm]{\texttt{127}}
        \item comment est codé le nombre entier relatif $-1$? \hfill\reponseLigne[2cm]{\texttt{11111111}}
    \end{itemize}

\end{exercice_cours}

\medskip

\end{document}